\section{Three-leaf K3P geometry}\label{sec:three-leaf}

Normalize \(q_{000}=1\) and order the sixteen conservation-supported
three-leaf coordinates as

\begin{equation}\label{eq:three-order}
\begin{split}
(&000,0CC,0GG,0TT,C0C,CC0,CGT,CTG,\\
 &G0G,GCT,GG0,GTC,T0T,TCG,TGC,TT0).
\end{split}
\end{equation}

The resulting affine chart has dimension \(15\).

\subsection{Trees and ordinary sunlets}

A three-leaf tree with pendant spectra \(a,b,c\) has

\[
 q_{xyz}=a_xb_yc_z\qquad(x+y+z=0).
\]

For the ordinary sunlet with reticulation adjacent to the third leaf, use
the exact map

\begin{equation}\label{eq:sunlet-map}
 q_{xyz}=a_xb_yc_z
 \bigl[\lambda f_y d_z+(1-\lambda)f_xe_z\bigr],
 \qquad x+y+z=0.
\end{equation}

The other orientations are obtained by labelled port transport.

For two triples \(L=(u,v,w)\) and \(R=(u',v',w')\), write
\(I(L;R)=q_uq_vq_w-q_{u'}q_{v'}q_{w'}\).  Define the six cubic circuits

\begin{align*}
I_1&=I((000,CGT,GTC);(0TT,C0C,GG0)),\\
I_2&=I((000,CTG,TGC);(0GG,C0C,TT0)),\\
I_3&=I((000,GCT,TGC);(0CC,GG0,T0T)),\\
I_4&=I((000,GTC,TCG);(0CC,G0G,TT0)),\\
I_5&=I((000,CTG,GCT);(0TT,CC0,G0G)),\\
I_6&=I((000,CGT,TCG);(0GG,CC0,T0T)).
\end{align*}

\begin{proposition}[Tree--ordinary-sunlet separation]\label{prop:tree-sunlet}
The observable

\[
 \mathcal S=\sum_{j=1}^6 I_j^2
\]

vanishes identically on the three-leaf K3P tree model and is strictly
positive on every ordinary three-sunlet with strict \(\Dplus\) edges and
\(0<\lambda<1\).  Leaf permutations cover all three orientations.
\end{proposition}

\begin{proof}
Tree monomials make every \(I_j\) vanish.  On \eqref{eq:sunlet-map}, each
pullback factors as a positive arm and inheritance monomial times two cross
factors involving one of

\[
 A_C=d_C-d_Gd_T,\qquad
 A_G=d_G-d_Cd_T,\qquad
 A_T=d_T-d_Cd_G.
\]

Sector by sector, if \(A_h\ne0\), vanishing of the two circuits carrying
that composition margin forces the two reciprocal cross ratios to equal
\(d_h\), hence \(d_h^2=1\), impossible for \(0<d_h<1\).  If all three
\(A_h\) vanish, their product gives
\(p=p^2\) for \(p=d_Cd_Gd_T\in(0,1)\), also impossible.  Therefore at least
one circuit is nonzero.  The six exact factorizations are printed in the
supplement and regenerated from the literal map.
\end{proof}

This separator is pointwise on the complete three-leaf maps; it does not
classify a pruned mixed graph by topology alone.

\subsection{The triangle hypersurface}

Define the eight-term quartic

\begin{align}\label{eq:H14}
F_{H_{14}}={}&
 q_{000}q_{CGT}q_{GTC}q_{TCG}
-q_{000}q_{CTG}q_{GCT}q_{TGC}\nonumber\\
&-q_{0CC}q_{CGT}q_{G0G}q_{TT0}
+q_{0CC}q_{CTG}q_{GG0}q_{T0T}\nonumber\\
&+q_{0GG}q_{C0C}q_{GCT}q_{TT0}
-q_{0GG}q_{CC0}q_{GTC}q_{T0T}\nonumber\\
&-q_{0TT}q_{C0C}q_{GG0}q_{TCG}
+q_{0TT}q_{CC0}q_{G0G}q_{TGC}.
\end{align}

Let \(H_{14}=Z(F_{H_{14}})\subset\mathbb A^{15}\).

\begin{theorem}[Ordinary-triangle geometry]\label{thm:H14}
The three labelled ordinary-triangle orientation maps have generic
normalized rank \(14\), have the same irreducible image closure \(H_{14}\),
and share a strict continuous-time smooth rank-\(14\) analytic germ.  They do
not contain an ambient-open \(15\)-dimensional germ.
\end{theorem}

\begin{proof}
Exact substitution of all three transported versions of
\eqref{eq:sunlet-map} annihilates \eqref{eq:H14}.  Give edges
\(a,b,c,d,e\) the isotropic spectrum \((1/2,1/2,1/2)\), give \(f\) the
spectrum \((1/3,1/3,1/3)\), and set \(\lambda=1/2\).  Every orientation then
maps to the same tensor

\[
q_0=\left(1,\frac1{12},\frac1{12},\frac1{12},
\frac1{12},\frac1{12},\frac1{48},\frac1{48},
\frac1{12},\frac1{48},\frac1{12},\frac1{48},
\frac1{12},\frac1{48},\frac1{48},\frac1{12}\right).
\]

All edges are strict continuous-time.  Exact \(14\)-minors at the three
preimages have determinants

\[
 -\frac1{760840571584512},\qquad
 -\frac1{760840571584512},\qquad
  \frac1{760840571584512}.
\]

The quartic gradient at \(q_0\) is nonzero; its six nonzero normalized
components have absolute value \(1/6912\).  Thus \(H_{14}\) is smooth and
has dimension \(14\) at \(q_0\).  Every orientation differential has rank
\(14\) and lies in the \(14\)-dimensional tangent kernel of \(dF_{H_{14}}\),
so it submerses onto a relative neighborhood in \(H_{14}\).  Intersect the
three neighborhoods and apply the constant-rank theorem to obtain physical
analytic sections.

For irreducibility, regard the normalized quartic as linear in \(q_{0CC}\).
Its coefficient is the primitive disjoint-support binomial

\[
 -q_{CGT}q_{G0G}q_{TT0}+q_{CTG}q_{GG0}q_{T0T}.
\]

The exponent difference is primitive, so this coefficient is irreducible.
At the certified specialization the coefficient vanishes while the remainder
equals \(-1\); hence it does not divide the remainder.  Gauss's lemma proves
that \(F_{H_{14}}\) is irreducible.  Since every orientation is contained in
this hypersurface and contains a relative neighborhood of its smooth locus,
each image closure is exactly \(H_{14}\).  None is ambient-open in
\(\mathbb A^{15}\).
\end{proof}

For comparison, the normalized three-leaf tree map has rank \(9\); one exact
minor is \(8192/28940625\).  The outer-class double-theta map in
\cref{sec:outer-obstruction} has rank \(15\).  The rank-\(14\) triangle
ambiguity is therefore intermediate and must be contextualized relative to
\(H_{14}\), not to the ambient tensor chart.

\begin{lemma}[Contextual \(H_{14}\) germ]\label{lem:H14-context}
Replace one ordinary triangle orientation by another inside an otherwise
unchanged labelled context, possibly reconnecting two triangle terminals
inside a theta factor.  The two complete network maps share a common germ
that is full-dimensional relative to both images.
\end{lemma}

\begin{proof}
Let \(U\subset H_{14}\) be the common relative germ from
\cref{thm:H14}, \(C\) a physical chart of the unchanged context, and write
the common labelled multilinear contraction as

\[
 \Psi(u,c)(g)=\sum_{\mathbf h\in\G^3}C(g;\mathbf h,c)u(\mathbf h).
\]

Choose one point of \(U\times C\) where \(D\Psi\) has its generic maximal
rank and shrink to a constant-rank neighborhood.  Physical sections of each
triangle orientation through \(U\) give the same lower rank for each
contextual map; direct factorization of each map through \(\Psi\) gives the
reverse inequality.  A constant-rank section of \(\Psi\), followed by the
orientation sections, yields one common full-dimensional contextual germ.
No tensor-product independence among the three terminals is assumed.
\end{proof}
